\chapter{Network Flows}

\section{Networks, flows, and cuts}

\begin{notedefinition}
A \emph{network} is a quadruple
\[
  N=(D,c,s,t),
\]
where $D=(V,A)$ is a finite directed graph, $c\colon A\to\R_{\ge0}$ is the
\emph{capacity function}, and $s,t\in V$ are distinct vertices called the
\emph{source} and the \emph{sink}.

For $X\subseteq V$ and a function $g\colon A\to\R$, put
\[
\begin{aligned}
  \delta^+(X)&:=\{xy\in A:x\in X,\ y\notin X\},\\
  \delta^-(X)&:=\{xy\in A:x\notin X,\ y\in X\},\\
  g_{\out}(X)&:=\sum_{e\in\delta^+(X)}g(e),
  &g_{\inflow}(X)&:=\sum_{e\in\delta^-(X)}g(e).
\end{aligned}
\]
For a single vertex $v$, we abbreviate $g_{\out}(\{v\})$ and
$g_{\inflow}(\{v\})$ to $g_{\out}(v)$ and $g_{\inflow}(v)$.
\end{notedefinition}

\begin{notedefinition}
A \emph{flow} in $N$ is a function $f\colon A\to\R_{\ge0}$ satisfying
\begin{enumerate}
  \item $0\le f(e)\le c(e)$ for every $e\in A$;
  \item $f_{\inflow}(v)=f_{\out}(v)$ for every $v\in V\setminus\{s,t\}$.
\end{enumerate}
The second condition is called \emph{flow conservation}.  The \emph{value} of
$f$ is
\[
  \val(f):=f_{\out}(s)-f_{\inflow}(s).
\]
The zero function is the \emph{zero flow}.
\end{notedefinition}

\begin{notedefinition}
A flow is \emph{maximum} if its value is at least the value of every other
flow in the same network.
\end{notedefinition}

\begin{notetheorem}
Let $f$ be a flow in $N=(D,c,s,t)$.  If $S\subseteq V$ satisfies
$s\in S$ and $t\notin S$, then
\[
  f_{\out}(S)-f_{\inflow}(S)=\val(f).
\]
In particular,
\[
  \val(f)=f_{\inflow}(t)-f_{\out}(t).
\]
\end{notetheorem}
\begin{proof}
Summing $f_{\out}(v)-f_{\inflow}(v)$ over $v\in S$ cancels every arc whose
two ends lie in $S$.  The uncancelled terms are exactly the arcs leaving or
entering $S$, so
\[
  f_{\out}(S)-f_{\inflow}(S)
  =\sum_{v\in S}\bigl(f_{\out}(v)-f_{\inflow}(v)\bigr).
\]
All summands vanish by flow conservation except the one at $s$, which is
$\val(f)$.  Taking $S=V\setminus\{t\}$ gives the second identity.
\end{proof}
\begin{notedefinition}
An $(s,t)$-\emph{cut} is a partition $V=S\dot\cup T$ with $s\in S$ and
$t\in T$.  We write the cut as $[S,T]$ and define its capacity by
\[
  c[S,T]:=c_{\out}(S)=\sum_{e\in\delta^+(S)}c(e).
\]
A cut of minimum capacity is a \emph{minimum cut}.
\end{notedefinition}

\begin{notetheorem}[weak duality]
For every flow $f$ and every $(s,t)$-cut $[S,T]$,
\[
  \val(f)\le c[S,T].
\]
\end{notetheorem}
\begin{proof}
By Theorem~5.1 and the capacity constraints,
\[
  \val(f)=f_{\out}(S)-f_{\inflow}(S)
  \le f_{\out}(S)
  \le c_{\out}(S)=c[S,T].
\]
\end{proof}
\section{Augmenting paths and residual networks}

\begin{notedefinition}
Let
\[
  P=(v_0,v_1,\dots,v_k),\qquad v_0=s,\quad v_k=t,
\]
be a path in the underlying undirected graph of $D$.  For each step
$v_{i-1}v_i$, call it
\begin{itemize}
  \item \emph{forward} if $v_{i-1}v_i\in A$, with residual capacity
        $c(v_{i-1}v_i)-f(v_{i-1}v_i)$;
  \item \emph{backward} if $v_iv_{i-1}\in A$, with residual capacity
        $f(v_iv_{i-1})$.
\end{itemize}
The residual capacity of $P$ is the minimum of these quantities.  If it is
positive, $P$ is an $f$-\emph{augmenting path}.  Augmenting by $\varepsilon$
means adding $\varepsilon$ to every forward arc of $P$ and subtracting
$\varepsilon$ from every backward arc.
\end{notedefinition}

\begin{noteremark}
An augmenting path records exactly where more flow can be pushed forward and
where existing flow can be cancelled and rerouted backward.
\end{noteremark}

\begin{notetheorem}
For a flow $f$ in a network $N$, the following are equivalent.
\begin{enumerate}
  \item $f$ is a maximum flow.
  \item There is no $f$-augmenting path from $s$ to $t$.
  \item There is an $(s,t)$-cut $[S,T]$ such that
        $c[S,T]=\val(f)$.
\end{enumerate}
\end{notetheorem}
\begin{proof}
If an augmenting path has residual capacity $\varepsilon>0$, augmenting along
it preserves the capacity constraints and flow conservation and increases the
flow value by $\varepsilon$.  Thus (1) implies (2).

Assume (2).  Form the residual digraph $D_f$: for each arc $xy\in A$, include
$xy$ when $c(xy)-f(xy)>0$, and include $yx$ when $f(xy)>0$.  Let $S$ be the
vertices reachable from $s$ in $D_f$, and let $T=V\setminus S$.  Since there
is no augmenting path, $t\in T$.  If $xy\in A$ has $x\in S$, $y\in T$, then
its forward residual capacity is zero, so $f(xy)=c(xy)$.  If $yx\in A$ has
$x\in S$, $y\in T$, then its backward residual capacity is zero, so
$f(yx)=0$.  Therefore
\[
\begin{aligned}
  \val(f)
  &=f_{\out}(S)-f_{\inflow}(S)\\
  &=\sum_{e\in\delta^+(S)}c(e)-0
   =c[S,T],
\end{aligned}
\]
and (3) holds.

Finally, (3) implies (1) by weak duality: every flow $g$ satisfies
$\val(g)\le c[S,T]=\val(f)$.
\end{proof}
\begin{notetheorem}[Max-Flow Min-Cut Theorem]
In every finite network, the maximum value of a flow equals the minimum
capacity of an $(s,t)$-cut.
\end{notetheorem}
\begin{proof}
A maximum flow exists because the feasible flows form a nonempty compact
polytope and the value is continuous.  Apply Theorem~5.3 to such a flow.
Weak duality gives the reverse inequality, so equality follows.
\end{proof}

\begin{notealgorithm}[Ford--Fulkerson]
Start with the zero flow.  While an augmenting path $P$ exists, augment by the
full residual capacity of $P$.  When no augmenting path remains, return the
current flow.
\end{notealgorithm}

\begin{notetheorem}
If every capacity is an integer, Ford--Fulkerson terminates after finitely many
augmentations and returns an integer-valued maximum flow.  In particular,
every network with integer capacities has an integer-valued maximum flow.
\end{notetheorem}
\begin{proof}
Starting from the zero flow, every residual capacity and every augmentation
amount is an integer.  Each augmentation increases the flow value by at least
$1$.  Weak duality bounds that value by the capacity of any fixed cut, so only
finitely many augmentations are possible.  At termination Theorem~5.3 gives a
maximum flow.
\end{proof}

\begin{noteremark}
For arbitrary real capacities, a poor choice of augmenting paths can lead to an
infinite sequence of augmentations.  The integer-capacity hypothesis is what
makes the elementary termination argument valid.
\end{noteremark}

\begin{notedefinition}
The \emph{residual digraph} $D_f$ of a flow $f$ has vertex set $V(D)$ and,
for every original arc $xy$, contains
\begin{itemize}
  \item a forward residual arc $xy$ of capacity $c(xy)-f(xy)$ when this is
        positive;
  \item a backward residual arc $yx$ of capacity $f(xy)$ when this is
        positive.
\end{itemize}
Parallel residual arcs may occur.  An augmenting path is precisely a directed
$s$--$t$ path in $D_f$.
\end{notedefinition}
\section{Applications}

\begin{noteapplication}[vertex-disjoint paths by vertex splitting]
Let $G$ be an undirected graph and let $u,v$ be distinct nonadjacent vertices.
Construct a network as follows.
\begin{enumerate}
  \item Keep $u$ as the source and $v$ as the sink.
  \item Replace every vertex $x\in V(G)\setminus\{u,v\}$ by two vertices
        $x^-$ and $x^+$ and the arc $x^-x^+$ of capacity $1$.
  \item For every edge $xy\in E(G)$, add arcs that represent traversal in
        both directions: from the output copy of one end to the input copy of
        the other.  Give these arcs a capacity larger than $|V(G)|$, or simply
        capacity $\infty$.
\end{enumerate}
Then the value of a maximum flow equals the maximum number of internally
vertex-disjoint $u$--$v$ paths.  By Max-Flow Min-Cut, it also equals the
minimum size of a $u$--$v$ vertex separator.
\end{noteapplication}
\begin{proof}
Every integral unit of flow determines a $u$--$v$ route.  The capacity-one arc
$x^-x^+$ permits at most one route to use an internal vertex $x$.  Conversely,
a family of internally disjoint paths gives an integral flow by sending one
unit along each corresponding route.  All finite minimum cuts can therefore
be chosen among the split arcs, and their capacities count deleted vertices.
\end{proof}

\begin{noteremark}
Taking the minimum over all relevant pairs $u,v$ recovers the vertex
connectivity $\kappa(G)$.  Adjacent pairs require the usual harmless
modification that forbids the direct edge from bypassing an internal
separator.
\end{noteremark}
\begin{noteapplication}[baseball elimination]
Suppose teams $1,\dots,n$ have current wins $w_i$, and let $r_{ij}$ be the
number of remaining games between teams $i$ and $j$.  To test whether team
$n$ can still finish tied for first, put
\[
  W:=w_n+\sum_{i<n}r_{in},
\]
its largest possible final number of wins.  If some $w_i>W$, team $n$ is
already eliminated.  Otherwise build a network with
\begin{itemize}
  \item a game vertex $x_{ij}$ for every $1\le i<j<n$;
  \item a team vertex $y_i$ for every $i<n$;
  \item arcs $s x_{ij}$ of capacity $r_{ij}$;
  \item arcs $x_{ij}y_i$ and $x_{ij}y_j$ of infinite capacity;
  \item arcs $y_i t$ of capacity $W-w_i$.
\end{itemize}
Team $n$ is not eliminated if and only if a maximum flow saturates every arc
leaving $s$.
\end{noteapplication}
\begin{proof}
A unit sent from $x_{ij}$ to $y_i$ assigns one remaining game between $i$ and
$j$ as a win for team $i$.  Saturating all source arcs assigns every remaining
game not involving team $n$, while the capacity $W-w_i$ prevents team $i$
from exceeding team $n$'s best possible total.  The correspondence is
reversible.
\end{proof}
\begin{notedefinition}[circulation with demands]
Let $D=(V,A)$ have capacities $c\colon A\to\R_{\ge0}$ and demands
$d\colon V\to\R$.  Interpret $d(v)>0$ as net demand and $d(v)<0$ as net
supply.  A \emph{circulation satisfying $d$} is a function
$f\colon A\to\R_{\ge0}$ such that
\[
  f(e)\le c(e)\quad(e\in A),
  \qquad
  f_{\inflow}(v)-f_{\out}(v)=d(v)\quad(v\in V).
\]
Necessarily $\sum_{v\in V}d(v)=0$.
\end{notedefinition}

\begin{noteapplication}[reducing circulation to maximum flow]
Assume $\sum_v d(v)=0$.  Add a super-source $s$ and a super-sink $t$.
For each supply vertex $v$ with $d(v)<0$, add $sv$ of capacity $-d(v)$.
For each demand vertex $v$ with $d(v)>0$, add $vt$ of capacity $d(v)$.
Retain all original arcs and capacities.  A feasible circulation exists if
and only if a maximum $s$--$t$ flow saturates every arc leaving $s$ (and hence
every arc entering $t$).
\end{noteapplication}
\begin{proof}
Given a feasible circulation, append the prescribed amounts on the new arcs.
The demand equation becomes ordinary flow conservation at every original
vertex, so the resulting flow saturates all source arcs.  Conversely, delete
the new arcs from any flow that saturates them.  Flow conservation then gives
$f_{\inflow}(v)-f_{\out}(v)=d(v)$ at each original vertex.
\end{proof}
